\documentclass[11pt]{article}
\usepackage{fontspec}
\setmainfont{Times New Roman}
\setsansfont{Arial}
\setmonofont{Consolas}
\usepackage{amsmath,amssymb}
\usepackage{geometry}
\usepackage{microtype}
\geometry{a4paper,margin=3cm}
\setlength{\parindent}{1.25em}
\setlength{\parskip}{0pt}
\emergencystretch=14em
\allowdisplaybreaks
\providecommand{\vdotswithin}[1]{\mathrel{\vdots}}

\begin{document}

\begin{center}
{\Large\bfseries 8. O cělom racionalnom predstavjenju invariantov sistema libovoljno mnogih osnovnyh form}\par
\vspace{1.2em}
Math. Ann. 77 (1916), s. 93--102
\end{center}

\vspace{1em}

Na koncu svojego prispevka k Švarcovomu jubilejnomu sborniku \emph{Über die Invarianten eines Systems von beliebig vielen Grundformen} D. Hilbert izkazuje predpoloženje, že te invarianty možno cělo i racionalno predstaviti črez konečno mnoge invarianty sistema $(J,PJ)$, gdě $J$ označuje polny sistem invariantov linearno nezavisnyh osnovnyh form, a invarianty $PJ$ izhodet iz $J$ črez polarne procesy.

V slědujučem pokazujem v I, že to predpoloženje stvarno jest pravilno, i to kako prosty poslědok redukcijnogo teorema: vsaku formu od libovoljno mnogih redov po $n$ varijablov v vsakom možno predstaviti kako sumu polarov form, ktore soderžet samo $n$ fiksovanyh redov i same sut izvedene iz izhodnoj formy polarnymi procesami. Toj redukcijny teorem jest specialny slučaj, prvy put formulovany F. Mertensom, obobčenja Klebš--Gordanovogo razloženja po redah na formy od libovoljno mnogih redov po $n$ varijablov; to obobčenje, dano A. Kapellijem, F. Mertensom i J. Derjutsom, dokazano jest formalnym prěobraženjem procesov diferencialovanja.\footnote{F. Mertens: \emph{Über eine Formel der Determinantentheorie}. (Formula 5) Sitzb. d. Ak. d. Wiss. Wien, Bd. 91, II. Abt. (1885), i podrobněje (s priměnjenjami) v: \emph{Über ganze Funktionen von $m$ Systemen von je $n$ Unbestimmten}. Monatsh. f. Math. u. Phys. 4. Jahrg. (1893). Za někotre dokazy Kapellija (od 1882) sr. napr.: Math. Ann. 37, ili avtografovane \emph{Lezioni sulla teoria delle forme algebriche} (1902). J. Deruyts: \emph{Essai d'une theorie générale des formes algébriques}, Mém. d. 1. Soc. Roy. des Sciences de Liège (2) 17 (1892).}

V II dajem ješče bolěje ponjatijny dokaz redukcijnogo teorema, razmatrajut čto vprašanje kako problem ekvivalentnosti linearnyh familij form. To gledišče, i takože suštstvenu čest upotrěbljenyh razsmatrjanj, berem iz raboty \emph{Über algebraische Modulsysteme}\footnote{E. Fischer: Über algebraische Modulsysteme und lineare homogene partielle Differentialgleichungen mit konstanten Koeffizienten, \S\ 1--4, J. f. M. 140 (1911).} i iz s njom svęzanyh nepublikovanyh teoremov E. Fišera, ktory mi ljubazno dozvolil jih upotrěbljenje.

Nakonec pokazujem v III, kako odgovarajuča razsmatrjanja vodet takože k najobčejšim razloženjam po redah.

\subsection*{I.}
Nehaj $\theta$ bude forma, homogena v vsakom iz $N>n$ redov
\[
  A_1,A_2,\ldots,A_N,
\]
gdě občno red $A_k$ može sostojati iz $n$ elementov
\[
  A_k^{(1)},A_k^{(2)},\ldots,A_k^{(n)}.
\]
Koeficienty $\theta$ mogut byti neopreděljene ili veličiny danej oblasti racionalnosti. Nadalje nehaj $P$ označuje cělu racionalnu funkciju -- s racionalnymi cělymi koeficientami -- od polarnyh procesov
\[
  P_{hk}=\left(A_h\frac{\partial}{\partial A_k}\right)
  =A_h^{(1)}\frac{\partial}{\partial A_k^{(1)}}+
   A_h^{(2)}\frac{\partial}{\partial A_k^{(2)}}+
   \cdots+
   A_h^{(n)}\frac{\partial}{\partial A_k^{(n)}},
\]
pri čem pod produktom $P_{hk}P_{ij}\theta$ treba razuměti priměnjenje $P_{hk}$ na $P_{ij}\theta$.

Togda spomenuty redukcijny teorem jest dan identičnostju
\begin{equation}
  \theta=\sum PZ,
\tag{1}
\end{equation}
gdě formy $Z$ soderžet ješče samo redy
\[
  A_1,A_2,\ldots,A_n
\]
i sut izvedene iz $\theta$ polarnymi procesami $P$.

Teper razsmatrjamo sistem osnovnyh form $(F')$, sostojajuči iz $N$ form jednakogo reda i jednakogo čisla varijablov
\[
  F_1,F_2,\ldots,F_N,
\]
ktorih koeficienty odgovarajuče berut se kako $N$ redov
\[
  A_1,A_2,\ldots,A_N.
\]
Nehaj $(J)$ označuje polny sistem -- sostojajuči iz konečno mnogih invariantov -- form $F_1,\ldots,F_n$, tako že vsaky invariant form $F_1,\ldots,F_n$ može se cělo i racionalno izraziti črez invarianty $(J)$. Hilbertovo predpoloženje, ktoro treba dokazati, tada jest, že za simultanne invarianty $S$ form $F_1,\ldots,F_N$ polny sistem jest dan konečno mnogimi invariantami $(J,PJ)$, gdě $PJ$ označuje vse invarianty, izvedime iz $J$ polarnymi procesami $P$.

V samom dělu, vsaky simultanny invariant $S$ sistema $(F')$ s racionalnymi cělymi koeficientami jest forma homogena v vsakom iz $N$ redov $A_1,\ldots,A_N$, s racionalnymi cělymi koeficientami; zato možno priměniti na nego identičnost (1). Poněže priměnjenje polarnyh procesov $P$ na $S$ poznato opět proizvodi invarianty, formy $Z$ sut takože invarianty, imenno simultanne invarianty form $F_1,\ldots,F_n$, poněže soderžet samo redy $A_1,\ldots,A_n$; zato po predpoloženju one sut cěle racionalne funkcije invariantov $(J)$. Taka identičnost (1) prěhodi v
\[
  S=\sum PY,
\]
gdě $Y$ označajut produkty stepenjev invariantov $(J)$. Ale stvarno izvršenje procesov $P$ na $Y$, po definiciji $P$, sostoji iz posledovateljnogo konečno-kratnogo izvršenja prostyh polarnyh operacij $P_{hk}$, ktore po pravilu diferencialovanja produktov vodet k sumam produktov iz invariantov $(J)$ i $(PJ)$. To samo važje za priměnjenje $P_{hk}$ na produkt iz $(J)$ i $(PJ)$; zato i $PY$ davaje samo cěle racionalne kombinacije iz $(J,PJ)$. {\itshape Těm jest dokazano cělo racionalno predstavjenje vseh invariantov $S$ črez $(J,PJ)$.}

\subsection*{II.}
Prěd dokazom redukcijnogo teorema prědpoložimo něčto o linearnyh familijah form. Pod {\itshape linearnoju familijeju form nad $K$} razumějemo sistem form jednakogo stepena, polny v tom smyslu, že zajedno s $\theta_1$ i $\theta_2$ on vsegda soderži takože $c_1\theta_1+c_2\theta_2$, gdě $c_1,c_2$ sut veličiny prědpisanej oblasti racionalnosti $K$, i gdě $K$ soderži oblast koeficientov form $\theta$. Medžu tymi formami vsegda jest konečno čislo -- recimo $\rho$ -- linearno nezavisnyh nad $K$, a vsake $\rho+1$ form zadovoljajut linearnu relaciju s koeficientami iz $K$; familija imaje {\itshape rang $\rho$} vzhodno k $K$.\footnote{Občejše linearne familije form -- ktoryh rang ne mora byti konečny -- dobivajut se, ako se opusti ograničenje jednakogo stepena, ili takože ograničenje, že $K$ soderži oblast koeficientov form $\theta$.} Upotrěbimo legko vidimy fakt: ako $\mathcal L$ i $\mathcal T$ sut dvě linearne familije nad $K$ jednakogo ranga $\rho$, i ako $\mathcal T$ jest podfamilija $\mathcal L$, togda $\mathcal L$ i $\mathcal T$ sut identične (ekvivalentne).

S tym prěhodimo k redukcijnomu teoremu. Identičnost iz I
\[
  \theta=\sum PZ
\]
prěhodi v analogičnu identičnost za $P\theta$, ako se na obě strany priměni ten samy polarny proces $P$; zato redukcijny teorem v to samo vrěme davaje predstavjenje vseh polarov $\theta$ sumami specialnyh polarov $PZ$. S drugoj strany, te specialne polary sigurno soderžet se medžu vsimi polarmi, takže redukcijny teorem tvrdi ekvivalentnost tyh dvoh linearnyh familij form; osoblivo takože ekvivalentnost tyh podfamilij, ktoryh formy v vsakom pojedinom redu imajut toj samy stepen kako $\theta$. Za dokaz toj ekvivalentnosti vstavimo ješče posrědnu familiju, opreděljenu formami $Z$. Zaradi prostoty najprvo dajemo dokaz v slučaju treh binarnyh redov $(N=3,n=2)$, a potom kratko naznačimo paralelny občy dokaz.

Nehaj zato $\theta(x,y,z)$ imaje odgovarajuče stepeny $\alpha,\beta,p$ v redah
\[
  x_1x_2;\quad y_1y_2;\quad z_1z_2;
\]
koeficienty $\theta$ najprvo nehaj budut neopreděljene $a$ (ili takože racionalne čisla). Tada razsmatrjamo tri linearne familije form.

1. Familija $\mathcal L$, sostojajuča iz vseh form $f(x,y,z)$, ktore nastali iz $\theta$ polarnymi procesami $P$ i ktore v pojedinyh redah $x,y,z$ imajut toj samy stepen kako $\theta$; $\mathcal L$ osoblivo soderži i $\theta$. Ako se $f(x,y,z)$ razsmatra kako forma od $x,y,z$ i neopreděljeń $a$, togda oblast koeficientov jest pole $R$ racionalnyh čisel; za $K$ takože treba vzeti $R$, poněže samo pri racionalno-cělyh $\theta_1,\theta_2$ izraz $c_1P_1+c_2P_2$ znovu stavaje procesom $P$; nehaj $\mathcal L$ imaje rang $\rho$ vzhodno k $R$.

2. Familija $\mathcal S$, sostojajuča iz vseh form $g(\lambda;x,y)$, opredělěnyh črez
\[
  g(\lambda;x,y)=f(x,y,\lambda_1x+\lambda_2y),
\]
ili, vyraženo polarnymi procesami,
\[
\begin{aligned}
  g(\lambda;x,y)
  &=\frac1{p!}\left[(\lambda_1x+\lambda_2y)\frac{\partial}{\partial z}\right]^p f(x,y,z) \\
  &=g_0(xy)\lambda_1^p+g_1(x,y)\lambda_1^{p-1}\lambda_2+\cdots+g_p(xy)\lambda_2^p,
\end{aligned}
\]
pri čem imajemo:
\[
  i!(p-i)!\,g_i(xy)=
  \left(y\frac{\partial}{\partial z}\right)^i
  \left(x\frac{\partial}{\partial z}\right)^{p-i}f(xyz).
\]
\footnote{To znači kako v I:
\[
 \left(t\frac{\partial}{\partial x}\right)=t_1\frac{\partial}{\partial x_1}+t_2\frac{\partial}{\partial x_2},
\]
zato osoblivo:
\[
 \left[(\lambda_1x+\lambda_2y)\frac{\partial}{\partial z}\right]
 = (\lambda_1x_1+\lambda_2y_1)\frac{\partial}{\partial z_1}
 + (\lambda_1x_2+\lambda_2y_2)\frac{\partial}{\partial z_2},
\]
\[
 \left[z\left(\frac{\partial}{\partial\lambda_1}\frac{\partial}{\partial x}
 +\frac{\partial}{\partial\lambda_2}\frac{\partial}{\partial y}\right)\right]
 =z_1\left(\frac{\partial}{\partial\lambda_1}\frac{\partial}{\partial x_1}
 +\frac{\partial}{\partial\lambda_2}\frac{\partial}{\partial y_1}\right)
 +z_2\left(\frac{\partial}{\partial\lambda_1}\frac{\partial}{\partial x_2}
 +\frac{\partial}{\partial\lambda_2}\frac{\partial}{\partial y_2}\right).
\]}
Tako $g_i(xy)$ davajut formy $Z$ identičnosti (1): $g$ sut racionalno-cěle formy v $x,y,\lambda$ i neopreděljenah $a$; nehaj $\mathcal S$ imaje rang $\sigma$ vzhodno k $R$.

3. Familija $\mathcal T$, ktora izhodi iz $\mathcal S$ črez obratny polarny proces, kako $\mathcal S$ izhodi iz $\mathcal L$; formy $h$ od $\mathcal T$ opreděljene sut črez
\[
\begin{aligned}
 h(x,y,z)
 &=\frac1{p!}\left[z\left(\frac{\partial}{\partial\lambda_1}\frac{\partial}{\partial x}
 +\frac{\partial}{\partial\lambda_2}\frac{\partial}{\partial y}\right)\right]^p g(\lambda;xy)\\
 &=h_0(xyz)+h_1(xyz)+\cdots+h_p(xyz),
\end{aligned}
\]
gdě po 2. važje:
\[
  h_i(xyz)=
  \left(z\frac{\partial}{\partial x}\right)^{p-i}
  \left(z\frac{\partial}{\partial y}\right)^i g_i(xy).
\]
Pojedine $h_i$, i zato takože $h$, sut formy $PZ$ i jih sumy; nehaj $\mathcal T$ imaje rang $\tau$ vzhodno k $R$.

Kako konstrukcija form $h(x,y,z)$ od $\mathcal T$ pokazyvaje, one izhodet iz $\theta$ polarnymi procesami $P$ i v redah $x,y,z$ pojedino imajut toj samy stepen kako $\theta$; familija $\mathcal T$ zato tvori podfamiliju $\mathcal L$, i po vyše spomenutom faktu za dokaz ekvivalentnosti ostavaje samo dokazati ravnost ranga. Pri tom dokazu specialna struktura $f(xyz)$ -- polary jednoj i toj samej formy $\theta$ -- dalje ne upotrěbljaje se.

{\itshape a) Za sravnenje ranga $\mathcal L$ i $\mathcal S$ služi pomocna lema a): iz $f(x,y,\lambda_1x+\lambda_2y)=0$ nužno slěduje $f(x,y,z)=0$.} To neposrědnje slěduje iz identične relacije medžu tremi binarnymi redami
\[
  (xy)z+(yz)x+(zx)y=0,
  \qquad (xy)=(x_1y_2-x_2y_1),\ldots,
\]
ktora imaje poslědok
\[
  f[x,y,(xy)x+(xz)y]=(xy)^p f(x,y,z).
\]
Iz $f(x,y,\lambda_1x+\lambda_2y)=0$ (identično v $\lambda$) zato, kako pokazuje substitucija $\lambda_1=(zy)$, $\lambda_2=(xz)$, i poněže $(xy)\ne0$, nužno slěduje $f(x,y,z)=0$.

Po toj lemi medžu formami iz $\mathcal S$ i iz $\mathcal L$ vlada jedno-jednoznačno odgovaranje. Iz
\[
  f_1(x,y,\lambda_1x+\lambda_2y)=g(\lambda;xy),\qquad
  f_2(x,y,\lambda_1x+\lambda_2y)=g(\lambda;xy)
\]
bo nužno slěduje
\[
  f_1(xyz)=f_2(xyz).
\]
Dalje, vsaka relacija v $\mathcal S$ sahranja se takože medžu jednoznačno odgovarajučimi formami v $\mathcal L$ (i obratno). Iz
\[
  c_1g^{(1)}(\lambda;xy)+c_2g^{(2)}(\lambda;xy)+\cdots+c_rg^{(r)}(\lambda;xy)=0
\]
bo po lemi nužno slěduje
\[
  c_1f^{(1)}(x,y,z)+c_2f^{(2)}(x,y,z)+\cdots+c_rf^{(r)}(x,y,z)=0.
\]
{\itshape Těm jest dokazano ravnost ranga $\mathcal L$ i $\mathcal S$; $\rho=\sigma$.}

{\itshape b) Za sravnenje ranga $\mathcal S$ i $\mathcal T$ odgovarajuče služi pomocna lema b): iz $h(x,y,z)=0$ nužno slěduje $g(\lambda;xy)=0$.} Za dokaz pomyslimo, že po 3. $h(x,y,z)=0$ ravnoznačno jest izčezanju pojedinyh produktov stepenjev
\[
 \left(\frac{\partial}{\partial\lambda_1}\frac{\partial}{\partial x_1}
 +\frac{\partial}{\partial\lambda_2}\frac{\partial}{\partial y_1}\right)^{p_1}
 \left(\frac{\partial}{\partial\lambda_1}\frac{\partial}{\partial x_2}
 +\frac{\partial}{\partial\lambda_2}\frac{\partial}{\partial y_2}\right)^{p_2}g(\lambda;x,y),
 \qquad p_1+p_2=p,
\]
iz čego daljnje diferencialovanje davaje takože izčezanje produktov stepenjev
\[
\begin{aligned}
&\left(\frac{\partial}{\partial x_1}\right)^{\alpha_1}
 \left(\frac{\partial}{\partial x_2}\right)^{\alpha_2}
 \left(\frac{\partial}{\partial y_1}\right)^{\beta_1}
 \left(\frac{\partial}{\partial y_2}\right)^{\beta_2} \\
&\quad\cdot
 \left(\frac{\partial}{\partial\lambda_1}\frac{\partial}{\partial x_1}
 +\frac{\partial}{\partial\lambda_2}\frac{\partial}{\partial y_1}\right)^{p_1}
 \left(\frac{\partial}{\partial\lambda_1}\frac{\partial}{\partial x_2}
 +\frac{\partial}{\partial\lambda_2}\frac{\partial}{\partial y_2}\right)^{p_2}
 g(\lambda;xy)
\end{aligned}
\]
Zato izčezaje takože vsaka linearna kombinacija tyh produktov stepenjev, i zato vsaka forma
\[
  f\left(\frac{\partial}{\partial x},\frac{\partial}{\partial y},
  \frac{\partial}{\partial\lambda_1}\frac{\partial}{\partial x}
  +\frac{\partial}{\partial\lambda_2}\frac{\partial}{\partial y}\right)g(\lambda,xy).
\]
Ako vyberemo $f$ specialno tako, že
\[
  f(x,y,\lambda_1x+\lambda_2y)=g(\lambda;xy),
\]
togda dobivamo takože
\[
  g\left(\frac{\partial}{\partial x};
    \frac{\partial}{\partial x},
    \frac{\partial}{\partial y}\right)g(\lambda,x,y)=0;
\]
ili -- ako položimo
\[
  g(\lambda,xy)=\sum G_i\lambda_1^{\nu_1}\lambda_2^{\nu_2}
  x_1^{\alpha_1}x_2^{\alpha_2}y_1^{\beta_1}y_2^{\beta_2},
\]
gdě $G_i$ sut racionalno-cěle linearne formy neopreděljeń $a$ -- 
\[
  \sum \nu_1!\nu_2!\alpha_1!\alpha_2!\beta_1!\beta_2!\,G_i^2=0
\]
i zato $g(\lambda;xy)=0$.

Iz pomocnoj lemy b), točno kako v a), slěduje {\itshape ravnost ranga $\mathcal S$ i $\mathcal T$; $\sigma=\tau$.}

Zato po a) i b) imajemo $\rho=\tau$; i těm jest -- poněže $\mathcal T$ jest podfamilija $\mathcal L$ -- dokazano ekvivalentnost tyh dvoh linearnyh familij. Vsaku formu iz $\mathcal L$, osoblivo takože $\theta$, možno zato predstaviti kako linearnu racionalno-cělu kombinaciju $\rho$ linearno nezavisnyh form $h(x,y,z)$:
\[
  \theta=\sum c_i h^{(i)}(x,y,z)=\sum PZ.
\]
Ta identičnost, izvedena za neopreděljene koeficienty $a$ formy $\theta$, ostavaje validna, ako neopreděljene $a$ zaměniti veličinami kakove-koli oblasti racionalnosti, i {\itshape redukcijny teorem jest tako občno dokazan v slučaju treh binarnyh redov.}

Občy dokaz za $N$ redov po $n$ varijablov teče sovsěm paralelno. Nehaj $\theta(A)$ imaje odgovarajuče stepen $\alpha_i$ v redu $A_i$. Tada znovu razsmatrjamo tri linearne familije form:

1. familiju $\mathcal L$, sostojajuču iz vseh form $f(A)$, ktore izhodet iz $\theta(A)$ polarnymi procesami $P$ i v vsakom pojedinom redu $A_i$ imajut toj samy stepen kako $\theta(A)$;

2. familiju $\mathcal S$, sostojajuču iz vseh form $g(\lambda;A_1\cdots A_n)$, ktore izhodet iz $f(A)$ črez substitucije
\[
\begin{aligned}
 A_{n+1}&=\lambda_{n+1}^{(1)}A_1+\lambda_{n+1}^{(2)}A_2+
        \cdots+\lambda_{n+1}^{(n)}A_n,\\
 &\vdotswithin{=}\\
 A_N&=\lambda_N^{(1)}A_1+\lambda_N^{(2)}A_2+
        \cdots+\lambda_N^{(n)}A_n
\end{aligned}
\]
pri tom koeficienty pojedinyh produktov stepenjev $\lambda$ v $g(\lambda;A_1\cdots A_n)$ davajut formy $Z$ identičnosti (1);

3. familiju $\mathcal T$, sostojajuču iz vseh form $h(A)$, opredělěnyh črez
\[
\begin{aligned}
 h(A)=&\left[A_{n+1}\left(
 \frac{\partial}{\partial\lambda_{n+1}^{(1)}}\frac{\partial}{\partial A_1}
 +\cdots+
 \frac{\partial}{\partial\lambda_{n+1}^{(n)}}\frac{\partial}{\partial A_n}\right)\right]^{\alpha_{n+1}} \\
&\cdots
\left[A_N\left(
 \frac{\partial}{\partial\lambda_N^{(1)}}\frac{\partial}{\partial A_1}
 +\cdots+
 \frac{\partial}{\partial\lambda_N^{(n)}}\frac{\partial}{\partial A_n}\right)\right]^{\alpha_N}
 g(\lambda;A_1\cdots A_n).
\end{aligned}
\]
$h(A)$ jest suma form $PZ$.

Teper, zaradi identične relacije medžu vsakymi $n+1$ redami, pomocna lema a) ostavaje validna; takože važje pomocna lema b), v dokazě ktorej čislo varijablov nikako ne vhodilo. $\mathcal L$ i $\mathcal T$ imajut zato jednaky rang i -- poněže $\mathcal T$ jest podfamilija $\mathcal L$ -- sut identične. Zato identičnost
\[
  \theta=\sum PZ
\]
važje za neopreděljene koeficienty, i slědovateljno za veličiny libovoljnoj oblasti racionalnosti kako koeficienty; {\itshape redukcijny teorem jest těm dokazany občno.}

\subsection*{III.}
Pokazujemo ješče kratko, kako analogične razsmatrjanja davajut takože {\itshape obče razloženje po redah} (za $N\le n$). Nehaj $Z$ bude pojedino homogena forma $n$ redov $A_1,\ldots,A_n$ po $n$ veličin; i nehaj $\Delta$ bude determinant tyh redov. Najprvo dokazujemo identičnost, odgovarajuču (1):
\begin{equation}
  Z=\sum PH \pmod{\Delta},
\tag{2}
\end{equation}
gdě formy $H$ soderžet ješče samo redy $A_1,\ldots,A_{n-1}$ i sut izvedene iz $Z$ procesami $P$.

Dokaz osnivaje se na prěobraženju relacij, izvedenyh v II, v kongruencije modulo $\Delta$. Zaradi prostoty neka osnovoju budut tri ternarne redy $x,y,z$; koeficienty predpoložene sut kako neopreděljene.

Tri linearne familije $\mathcal L,\mathcal S,\mathcal T$ nehaj budut opreděljene kako v II v slučaju binarnyh redov; jih rang nehaj bude $\rho,\sigma,\tau$, a $\rho^*$ i $\tau^*$ označajut rang $\mathcal L$ i $\mathcal T$ modulo $\Delta$ (t. j. jest $\rho^*$, ale ne $\rho^*+1$ form iz $\mathcal L$, ktore sut linearno nezavisne modulo $\Delta$). Tada važje kako v II

{\itshape Pomocna lema a): iz $f(x,y,\lambda_1x+\lambda_2y)=0$ nužno slěduje $f(xyz)\equiv0\pmod{\Delta}$ (i obratno).} Imenno po relaciji medžu četyrmi ternarnymi redami
\[
  \Delta_1x-\Delta_2y+\Delta_3z=\Delta t,
  \qquad \Delta_1=(yzt),\ldots,
\]
vsegda jest medžu tremi ternarnymi redami linearna relacija modulo $\Delta$:
\[
  \Delta_1x-\Delta_2y+\Delta_3z\equiv0\pmod{\Delta},
\]
i zato, kako v II, imajemo
\[
  f(x,y,-\Delta_1x+\Delta_2y)=\Delta_3^p f(x,y,z)\pmod{\Delta}.
\]
Iz $f(x,y,\lambda_1x+\lambda_2y)=0$ (identično v $\lambda$) zato -- poněže $\Delta$ jest nerazložimy i $\Delta_3$ ne jest děljivy črez $\Delta$ -- nužno slěduje $f(x,y,z)\equiv0\pmod{\Delta}$. Obratno jest jasno iz $\Delta(x,y,\lambda_1x+\lambda_2y)=0$. Iz togo zaključajemo kako v II: $\rho^*=\sigma$.

Pomocna lema b) ostavaje s dokazom validna, i imajemo zato: $\sigma=\tau$.

Za dokaz $\tau=\tau^*$ služi

{\itshape Pomocna lema c): iz $h(x,y,z)\equiv0\pmod{\Delta}$ nužno slěduje $h(x,y,z)=0$.}
$h(x,y,z)$ bo, kako polara, izčezaje pri priměnjenju poznatogo $\Omega$-procesa; to se izražaje diferencialnym uravnanjem
\[
  \Delta\!\left(\frac{\partial}{\partial x},\frac{\partial}{\partial y},\frac{\partial}{\partial z}\right)h(x,y,z)=0.
\]
Ako imajemo $h(x,y,z)=\Delta(x,y,z)\cdot k(x,y,z)$, togda daljnim diferencialovanjem izčezaje takože
\[
  k\!\left(\frac{\partial}{\partial x},\frac{\partial}{\partial y},\frac{\partial}{\partial z}\right)
  \Delta\!\left(\frac{\partial}{\partial x},\frac{\partial}{\partial y},\frac{\partial}{\partial z}\right)
  =
  h\!\left(\frac{\partial}{\partial x},\frac{\partial}{\partial y},\frac{\partial}{\partial z}\right)h(x,y,z),
\]
i zato $h(x,y,z)$. Těm imajemo $\rho^*=\sigma=\tau=\tau^*$, i poněže $\mathcal T$ jest podfamilija $\mathcal L$, identičnost (2) jest dokazana; odgovarajuče razsmatrjanje davaje dokaz pri $n$ varijablah.\footnote{Po 3. imaje se odgovarajuča, črez determinantny faktor izčezajuča relacija za $\Omega(h)$; v operatorskom zapisu to jest uravnanje, nastajuče iz produktov
\[
 \left(\lambda_1\frac{\partial}{\partial \xi}+\lambda_2\frac{\partial}{\partial \eta}\right),
 \qquad
 \left[z\left(\frac{\partial}{\partial\lambda_1}\frac{\partial}{\partial \xi}
 +\frac{\partial}{\partial\lambda_2}\frac{\partial}{\partial \eta}\right)\right]
\]
ono izčezaje zaradi izčezanja determinantnogo faktora.}

Priměnjajuči identičnost (2) na množitelj $\Delta$, dobivaje se razloženje:
\[
  Z=\varphi_0+\Delta\varphi_1+\Delta^2\varphi_2+\cdots+\Delta^\mu\varphi_\mu,
\]
gdě $\varphi_i$ pojedino izčezajut pri priměnjenju $\Omega$-procesa. $i$-kratno povtorjenje procesa zato davaje -- poněže občno imajemo: $\Omega(\Delta\cdot\psi)=c_1\psi+c_2\Delta\cdot\Omega(\psi)$ --:
\[
  c\cdot \Omega^i(Z)\equiv\varphi_i\pmod{\Delta};
\]
s drugoj strany po (2) imajemo:
\[
  c\cdot\Omega^i(Z)\equiv\sum PH_i\pmod{\Delta},
\]
gdě $H_i$ takože izvedene sut iz formy $\Omega^i(Z)$, kako $H$ iz $Z$. Po pomocnoj lemi c) (v razširenju na $n$ varijablov) dobivaje se zato: $\varphi_i=\sum PH_i$, i slědovateljno:
\begin{equation}
  Z=\sum PH+\Delta\cdot\sum PH_1+\Delta^2\cdot\sum PH_2+
  \cdots+\Delta^\mu\cdot\sum PH_\mu
\tag{3}
\end{equation}
{\itshape kako najobčejše razloženje formy od $n$ redov po $n$ varijablov po polarah form samo od $n-1$ redov.}

Razloženje po redah za formy od $N$ redov po $n$ varijablov $(N<n)$ dobivaje se iz togo, kako poznato, prostoju substitucijeju. Položimo (za $N=3$, $n>3$):
\[
  u=\xi_1x+\xi_2y+\xi_3z,
  \quad v=\eta_1x+\eta_2y+\eta_3z,
  \quad w=\zeta_1x+\zeta_2y+\zeta_3z
\]
i razložimo $H(u,v,w)=Z(\xi,\eta,\zeta)$ po (3). Tada, poněže $\xi,\eta,\zeta$ nastupajut samo v svęzkah $u,v,w$, imajemo:
\[
  \left(\eta\frac{\partial}{\partial \xi}\right)=
  \left(v\frac{\partial}{\partial u}\right)\cdots,
  \qquad
  \Omega=\sum_{ikl}\left(\frac{\partial}{\partial u}
  \frac{\partial}{\partial v}
  \frac{\partial}{\partial w}\right)_{ikl}(xyz)_{ikl}
  =\nabla_{uvw}(xyz).
\]
Po specializaciji $\xi_1=\eta_2=\zeta_3=1$; $\xi_2=\xi_3=\cdots=\zeta_2=0$ forma $H(u,v,w)$ prěhodi v $H(x,y,z)$, $\nabla_{uvw}(xyz)$ do čislovogo faktora v $\nabla_{xyz}(xyz)=\nabla_{xyz}$, poněže priměnjenje $(y\partial/\partial x)$ itd. na determinanty $(xyz)_{ikl}$ čini jih izčeznuti. Poněže odgovarajuče važje takože pri $N$ redah, iz (3) slěduje razloženje:
\begin{equation}
  H=\sum P\Phi+\sum P\Phi_1+\cdots+\sum P\Phi_\mu,
\tag{4}
\end{equation}
gdě $\Phi_i$ izhodet iz $\nabla^i_{A_1\ldots A_N}H$ črez procesy $P$ i soderžet ješče samo redy $A_1,\ldots,A_{N-1}$, osim determinantnyh svęzkov nastupajučih v $\nabla^i$, ktore, kako spomenuto, pri polarizaciji ne vhodet v račun.

Ako zato zaměniti te determinanty v $\nabla^i$ neopreděljenymi $p$, to nastajuče formy možno znovu razložiti po (4); i tako nakonec dobivaje se razloženje
\begin{equation}
  H=\left[\sum P\psi(p)\right]_{p=\text{determinanty }A_1\cdots A_N},
\tag{5}
\end{equation}
gdě $\psi$ izhodet iz izhodnoj formy črez procesy $P$ i $\nabla_{A_1\ldots A_N}(p)$, $\nabla_{A_1\ldots A_{N-1}}(p)$ itd. Tymi razloženjami i jih kombinacijami -- napr. ako se na formy $Z$ nastupajuče v (1) priměni (3), a potom (4) i (5) -- izčerpane sut vse poznate razloženja za formy od libovoljno mnogih redov po $n$ kogredientnyh varijablov.

Erlangen, 5. januar 1915.

\end{document}
